% id: 81-11
% book: 베이직쎈 미적분1 · 05 도함수의 활용 (2)
% section: 유형 051 도함수의 그래프와 함수의 극대·극소
% figure: tikz:fig-81-11
% answer: 
% answer_source: 
% variant_level: 0 (원본 전사)
% 이 파일은 build.mjs 가 items.json 에서 만든다. 고칠 때는 items.json 을 고친다.
\smash{\raisebox{1.5mm}{\dmkichul{베이직쎈 미적분1 81쪽 11번}}}
\noindent\begin{minipage}[t]{0.58\linewidth}\vspace{0pt}\raggedright 사차함수 $f(x)$의 도함수 $y=f'(x)$의 그래프가 오른쪽 그림과 같다. $f(0)=2$일 때, 옳은 것만을 \textbf{보기}에서 있는 대로 고른 것은? \end{minipage}\hfill\begin{minipage}[t]{0.4\linewidth}\vspace{0pt}\centering \resizebox{\ifdim\width>\linewidth\linewidth\else\width\fi}{!}{% 81-11(81쪽) — 사차함수 $f(x)$ 의 도함수 $y=f'(x)$ 의 그래프($x=-2,\ 0,\ 2$ 에서 $x$축과 만난다 · $x=-1$ 에서 극대, $x=1$ 에서 극소). 스캔 화질이 낮아 TikZ 로 다시 그림.
% 원문에 있는 것만 옮긴다: 곡선 · $x=-1,\ 1$ 의 안내 점선 · 눈금 -2, -1, 1, 2 · 라벨 O, x, y, y=f'(x). 원문에 없는 함숫값은 적지 않는다.
\begin{tikzpicture}[x=0.52cm, y=0.50cm, line width=0.6pt]
  \draw[->, >=stealth, line width=0.7pt] (-2.85,0) -- (3.05,0) node[below=1pt, font=\small] {$x$};
  \draw[->, >=stealth, line width=0.7pt] (0,-1.55) -- (0,2.60) node[left=1pt, font=\small] {$y$};
  \draw[densely dashed, line width=0.4pt] (-1,0) -- (-1,1.75);
  \draw[densely dashed, line width=0.4pt] (1,0) -- (1,-1.05);
  \draw plot[smooth, tension=0.7] coordinates {(-2.4,-1.6) (-2.2,-0.85) (-2,0) (-1.7,1.0) (-1.4,1.6) (-1,1.75) (-0.7,1.6) (-0.4,1.1) (-0.15,0.45) (0,0) (0.25,-0.56) (0.6,-0.95) (1,-1.05) (1.4,-0.8) (1.7,-0.42) (2,0) (2.2,0.58) (2.45,1.3) (2.6,1.95)};
  \node[above left=-2pt and -4pt, font=\small] at (-2,0) {$-2$};
  \node[below=1pt, font=\small] at (-1,0) {$-1$};
  \node[below left=-2pt and -3pt, font=\small] at (0,0) {$\pt{O}$};
  \node[above=1pt, font=\small] at (1,0) {$1$};
  \node[below=1pt, font=\small] at (2,0) {$2$};
  \node[anchor=south, font=\small] at (2.05,2.20) {$y=f'(x)$};
\end{tikzpicture}
}\end{minipage}\par
\noindent \bogi{\begin{minipage}{0.96\linewidth}\raggedright ㄱ. $f(x)$는 $x=-2$, $x=2$에서 극솟값을 갖는다.\\[1mm] ㄴ. $f(-1)<f\left(-\dfrac{1}{2}\right)$\\[1mm] ㄷ. $f(x)$의 극댓값은 $2$이다.\end{minipage}}
\dmchoicesauto{}{ㄱ}{ㄷ}{ㄱ, ㄴ}{ㄱ, ㄷ}{ㄱ, ㄴ, ㄷ}
